% Shared preamble for the public math-only manuscript.
% Included by docs/main.tex via % Shared preamble for the public math-only manuscript.
% Included by docs/main.tex via % Shared preamble for the public math-only manuscript.
% Included by docs/main.tex via \input{preamble} as its first line.

\documentclass[11pt]{article}

\usepackage[margin=1in]{geometry}
\usepackage{amsmath}
\usepackage{amssymb}
\usepackage{amsfonts}
\usepackage{booktabs}
\usepackage{array}
\usepackage[hidelinks]{hyperref}

% Return-type notation used throughout the manuscript.
%   \Rv  value-gap return   (V_i - P_0)/P_0, horizon H_value
%   \Rp  realized price return,              horizon H_event
% These are kept here deliberately: conflating the two was the root
% error in the v1 model, and a shared definition keeps them distinct.
\newcommand{\Rv}{R^{\mathrm{val}}}
\newcommand{\Rp}{R^{\mathrm{px}}}
\newcommand{\E}{\mathbb{E}}
\newcommand{\Prob}{\mathbb{P}}
\newcommand{\Var}{\operatorname{Var}}
\newcommand{\se}{\operatorname{se}}

\author{Blake Paulson}
\date{September 2026}
 as its first line.

\documentclass[11pt]{article}

\usepackage[margin=1in]{geometry}
\usepackage{amsmath}
\usepackage{amssymb}
\usepackage{amsfonts}
\usepackage{booktabs}
\usepackage{array}
\usepackage[hidelinks]{hyperref}

% Return-type notation used throughout the manuscript.
%   \Rv  value-gap return   (V_i - P_0)/P_0, horizon H_value
%   \Rp  realized price return,              horizon H_event
% These are kept here deliberately: conflating the two was the root
% error in the v1 model, and a shared definition keeps them distinct.
\newcommand{\Rv}{R^{\mathrm{val}}}
\newcommand{\Rp}{R^{\mathrm{px}}}
\newcommand{\E}{\mathbb{E}}
\newcommand{\Prob}{\mathbb{P}}
\newcommand{\Var}{\operatorname{Var}}
\newcommand{\se}{\operatorname{se}}

\author{Blake Paulson}
\date{September 2026}
 as its first line.

\documentclass[11pt]{article}

\usepackage[margin=1in]{geometry}
\usepackage{amsmath}
\usepackage{amssymb}
\usepackage{amsfonts}
\usepackage{booktabs}
\usepackage{array}
\usepackage[hidelinks]{hyperref}

% Return-type notation used throughout the manuscript.
%   \Rv  value-gap return   (V_i - P_0)/P_0, horizon H_value
%   \Rp  realized price return,              horizon H_event
% These are kept here deliberately: conflating the two was the root
% error in the v1 model, and a shared definition keeps them distinct.
\newcommand{\Rv}{R^{\mathrm{val}}}
\newcommand{\Rp}{R^{\mathrm{px}}}
\newcommand{\E}{\mathbb{E}}
\newcommand{\Prob}{\mathbb{P}}
\newcommand{\Var}{\operatorname{Var}}
\newcommand{\se}{\operatorname{se}}

\author{Blake Paulson}
\date{September 2026}


\title{Pre-Earnings Valuation and Event Risk\\
  \large A Theoretical Framework for Fair Value, Implied Movement, and Scenario Uncertainty}

\begin{document}

\maketitle

\begin{abstract}
\noindent
This paper develops a theoretical framework for analyzing expected price
movement around scheduled earnings events. The framework separates four objects
that are often collapsed into one: fair value, the market-implied scale of the
event move, the probability of discrete reaction states, and the extent to
which a valuation gap can be incorporated into price over the event horizon.

The paper makes four principal contributions. First, it treats the at-the-money
straddle as a risk-neutral expectation of an absolute move rather than as the
size of a tail outcome. Second, it shows that a neutral state defined as an
interval absorbs part of that absolute move, producing a closed-form leakage
correction. Third, it derives the expected-value and risk expressions for both
symmetric and asymmetric three-state reaction models. In the asymmetric model,
the sign of expected value changes when the ratio of downside to upside move
magnitudes equals the ratio of bullish to bearish probabilities. Fourth, it
separates fair value from event price through an explicit convergence
coefficient and derives a sampling-error floor for the directional probability
tilt.

The paper is a theoretical and methodological specification. It does not
present fitted parameters, a backtest, or evidence that a deployable forecasting
system produces economic profits. Its central conclusion is therefore about
identification: an apparent pre-earnings edge is not determined by a valuation
gap alone. It depends jointly on event magnitude, directional probabilities,
downside asymmetry, convergence, and estimation uncertainty.

This paper does not produce a security-specific expected-value forecast. Nor
should it be treated as a model for producing alpha, a trading signal, or a
deployable investment recommendation. The equations define theoretical
identities and identify the inputs and empirical evidence that a future
forecasting study would need; they do not supply those inputs or establish that
an alpha exists.
\end{abstract}

\tableofcontents
\newpage

%======================================================================
\section{Introduction}
%======================================================================

Pre-earnings valuation models commonly combine a fair-value estimate with an
expected move and then report a single expected return. That compression hides
several distinct questions. A valuation model concerns the level at which a
security might be fairly priced. An options market supplies a price for
dispersion over a specified horizon. A scenario model assigns probabilities to
possible reactions. A short-horizon price response need not equal the change in
fair value. Finally, all of these quantities are estimated with error.

This paper develops a common language for those questions. It is deliberately
more cautious than a trading rule. The object is a coherent theoretical model
whose assumptions, identities, and failure modes are visible before any
empirical estimation is attempted.

\subsection{Research question}

Let $P_0$ be the price immediately before the event and let $P_1$ be the price
at the end of a fixed event window. The central question is:

\begin{quote}
Under what assumptions can a pre-earnings fair-value estimate imply a
statistically distinguishable expected event return, once market-implied
dispersion and asymmetric reaction risk are accounted for?
\end{quote}

The answer is not a single formula. It is a sequence of distinctions. The
framework first constructs a valuation gap, then describes event-scale movement,
then assigns scenario probabilities, and finally states how the valuation gap
can enter the event-price return. The layers can be estimated separately, but
they must not be silently added together when they describe the same price
movement.

\subsection{Theoretical contribution}

The paper contributes the following results.

\begin{enumerate}
\item \textbf{Event-scale identification.} An at-the-money straddle prices a
      risk-neutral mean absolute move, not a one-standard-deviation move and
      not the magnitude of a tail state. The three-state magnitude must
      therefore be matched to a probability-weighted absolute-move moment.
\item \textbf{Neutral-band correction.} When the neutral state is an interval,
      the interval contributes to the mean absolute move. Under a normal
      reference distribution, the fraction is
      $\ell_{\mathcal N}(p_N^{\mathrm{ref}})=1-\exp(-z^2/2)$, where
      $z=\Phi^{-1}((1+p_N^{\mathrm{ref}})/2)$.
\item \textbf{Asymmetric reaction result.} If downside and upside magnitudes
      differ, the expected-value sign condition is
      $p_U=\rho p_B$, where $\rho=m_B/m_U$. Equal tail probabilities are not
      sufficient for zero expected value once skew is admitted.
\item \textbf{Identification and precision.} A valuation gap reaches the event
      price through a convergence coefficient $\kappa$, while the directional
      edge is driven by a probability tilt whose sampling error must be compared
      with the tilt at the event counts available for the intended application.
\end{enumerate}

\subsection{Scope and non-claims}

This is a theoretical preprint, not an empirical finance study. It does not
claim that any parameter has been estimated from a completed event panel, that
the probability model is calibrated out of sample, or that a position-sizing
rule is deployable. Numerical values are explicitly labelled as identities,
assumptions, or illustrations. The paper does not disclose production
architecture, implementation details, private data workflows, or proprietary
forecasting components.

The framework is useful even before estimation because it identifies what an
empirical study would have to measure. A theoretical result that an edge changes
sign with an unidentified parameter is itself a result about what cannot be
inferred from the available inputs.

\subsection{Organization}

Section~\ref{sec:setup} defines the horizons, state space, and notation.
Section~\ref{sec:value} describes the fair-value layer. Section~\ref{sec:move}
derives the event-scale correction. Section~\ref{sec:prob} formulates scenario
probabilities and the market-implied scale benchmark. Section~\ref{sec:kappa} gives the
fair-value-to-price bridge. Section~\ref{sec:ev} combines the layers without
double counting and derives the principal closed forms. Section~\ref{sec:error}
studies uncertainty. Section~\ref{sec:implications} states empirical
implications and limitations. Appendices collect secondary decision material
and a fully labelled symbolic illustration.

%======================================================================
\section{Setup, Horizons, and Notation}
\label{sec:setup}
%======================================================================

\subsection{Two horizons}

An earnings model contains at least two horizons:

\[
\boxed{
\begin{aligned}
  H_{\mathrm{event}} &= \text{a prespecified earnings-event window},\\
  H_{\mathrm{value}} &= \text{the horizon over which fair value is assessed}.
\end{aligned}}
\]

A discounted valuation over $H_{\mathrm{value}}$ and an options-implied move
over $H_{\mathrm{event}}$ are different objects. They may be related only by an
explicit price-adjustment model. In particular, a multi-period valuation gap must
not be treated as though it were a one-day expected return.

\subsection{Price, value, and returns}

Let $P_0$ and $P_1$ denote the pre-event and post-event market prices. Let
$V_0$ and $V_1$ denote the corresponding model-implied fair values. Define

\[
\Rv_i=\frac{V_i-P_0}{P_0},
\qquad
\Rp_i=\frac{P_i-P_0}{P_0},
\]

where the subscript $i$ denotes a reaction state. The superscripts distinguish
the valuation return from the realized price return; they are not interchangeable.

\subsection{State space}

The event reaction is represented by an ordered state variable
$Y\in\{B,N,U\}$: bearish, neutral, and bullish. Let

\[
\Prob(Y=i)=p_i,
\qquad
p_i\ge 0,
\qquad
p_B+p_N+p_U=1.
\]

The three-state representation is a deliberate low-dimensional modeling choice.
It does not claim that the true return distribution has only three points or
that its tails are thin. A fixed threshold $a>0$ defines the states before the
outcomes are observed:

\[
Y=
\begin{cases}
B, & \Rp<-a,\\
N, & |\Rp|\le a,\\
U, & \Rp>a.
\end{cases}
\]

The threshold is part of the model definition and is not identified by the
theory. Choosing it after observing which partition produces a desirable
expected value would make the result a tuning exercise rather than a theoretical
prediction.

\subsection{Notation}

\begin{center}
\begin{tabular}{@{}lll@{}}
\toprule
Symbol & Meaning & Definition or restriction \\
\midrule
$P_0,P_1$ & pre- and post-event market price & \\
$V_0,V_1$ & pre- and post-event fair value & model-implied \\
$\Rv_i$ & value-gap return in state $i$ & $(V_i-P_0)/P_0$ \\
$\Rp_i$ & realized event-price return in state $i$ & $(P_i-P_0)/P_0$ \\
$Y$ & reaction state & $Y\in\{B,N,U\}$ \\
$p_B,p_N,p_U$ & state probabilities & sum to one \\
$a$ & neutral-state threshold & fixed before outcomes \\
$m_{IV}$ & straddle-implied absolute move & approximately $\E^{\mathbb Q}|\Rp|$ \\
$M$ & physical tail absolute-move budget & defined in Section~\ref{sec:move} \\
$m_B,m_U$ & downside and upside state magnitudes & positive \\
$\rho$ & downside-to-upside magnitude ratio & $m_B/m_U$ \\
$\kappa$ & value-to-price convergence coefficient & free estimand in the core model \\
$EV$ & expected event return & statistical expected value only \\
$\gamma$ & relative risk-aversion coefficient & utility parameter \\
$\Theta$ & enterprise value & never denoted by $EV$ \\
\bottomrule
\end{tabular}
\end{center}

\subsection{Identification status of inputs}

The equations below do not assign numerical values to the threshold $a$, the
forecast horizon, the physical adjustment $h$, the tail-asymmetry ratio $\rho$,
the convergence coefficient $\kappa$, the probability vector, the regularization
weight $\eta$, or the utility and cost parameters. They are inputs, restrictions,
or estimands for a later empirical design. The state count and the choice of a
reference distribution are modeling choices rather than empirical facts. The
unit scale in $\operatorname{Logistic}(0,1)$ is an identification normalization,
and constants such as $\sqrt{2/\pi}$ are derived only under the explicitly stated
normal reference model. Algebraic indices, probability normalization, and
indicator values are notation, not calibrations.

%======================================================================
\section{The Fair-Value Layer}
\label{sec:value}
%======================================================================

\subsection{Valuation as a range, not a false point estimate}

Different valuation methods often share earnings, growth, discount-rate, and
capital-structure inputs. Their errors are therefore correlated, so mechanically
averaging them does not create independent information. A conservative summary
is a range:

\[
\boxed{
V_0\in[V_{\min},V_{\max}],
\qquad
V_{\min}=\min_j V_j,
\quad
V_{\max}=\max_j V_j.
}
\]

When a single reference value is needed for exposition, let $\hat V_0$ be the
value selected by the stated valuation protocol and let $\varsigma_V$ be its
uncertainty or dispersion summary. Neither choice is identified by the equations:

\[
\hat V_0\in[V_{\min},V_{\max}],
\qquad
\varsigma_V\ge 0.
\]

Here $\varsigma_V$ is a protocol-defined uncertainty summary, not automatically
a sampling standard error. A point estimate becomes a meaningful empirical
object only when the valuation methods, their dependence, and their out-of-sample
performance are specified.

\subsection{Forward earnings valuation}

For a $k$-period forecast path,

\[
EPS_{t+k}=EPS_t\prod_{j=1}^{k}(1+g_{E,j}),
\qquad
V_{\mathrm{PE}}=\frac{EPS_{t+k}PE_{t+k}}{(1+r)^k}.
\]

The justified multiple should respect the curvature of a growth model. One
optional structural anchor is the Gordon-style expression

\[
PE^*=\frac{1-b}{r-g_E},
\qquad g_E<r,
\]

so a residual log-linear specification can be written as

\[
\ln PE_f
=\ln\left(\frac{1-b}{r-g_E}\right)
 +\alpha_1 ROIC+\alpha_2 M-\alpha_3\sigma-\alpha_4 C+\varepsilon.
\]

The anchor prevents a negative multiple and makes clear which terms are
theoretical and which would require estimation.

\subsection{Discounted cash flow}

For completeness, an enterprise-value DCF is

\[
\Theta_{\mathrm{DCF}}
=\sum_{t=1}^{T}\frac{FCF_t}{(1+WACC)^t}
 +\frac{TV_T}{(1+WACC)^T},
\qquad
TV_T=\frac{FCF_T(1+g_T)}{WACC-g_T},
\]

and the implied equity value per share is

\[
V_{\mathrm{DCF}}
=\frac{\Theta_{\mathrm{DCF}}+\mathit{Cash}-\mathit{Debt}}{n_{\mathrm{sh}}},
\qquad g_T<WACC.
\]

Terminal value is often the dominant term. The framework therefore treats the
width of the valuation range as part of the result rather than hiding it inside
a precise-looking discounted cash-flow output.

\subsection{Valuation gap}

The pre-event value gap is

\[
\boxed{
D_V=\frac{V_0-P_0}{P_0}.
}
\]

Positive $D_V$ means that the model's fair value exceeds the current price. If
$\hat V_0$ and $\varsigma_V$ are used,

\[
D_V\pm\frac{\varsigma_V}{P_0}
\]

is more honest than a point estimate. A gap smaller than the valuation
dispersion is not evidence of detectable mispricing.

\subsection{Persistent fundamental regimes}

If a valuation path depends on a latent growth regime $G_t$, Bayesian updating
belongs to that regime rather than to the already-realized event outcome. Let
$G_t\in\{G_B,G_N,G_U\}$. Then

\[
\Prob(G_t\mid D_t)
=\frac{\Prob(D_t\mid G_t)\Prob(G_t)}
{\sum_j\Prob(D_t\mid G_j)\Prob(G_j)},
\]

and persistence can be represented by

\[
A_{ij}=\Prob(G_{t+1}=j\mid G_t=i),
\qquad
\Prob(G_{t+1}=j)=\sum_i\Prob(G_t=i\mid D_t)A_{ij}.
\]

This is a structural option for the valuation layer. It does not provide an
empirical estimate of the transition matrix.

%======================================================================
\section{The Event-Scale Layer}
\label{sec:move}
%======================================================================

\subsection{What an at-the-money straddle measures}

For an at-the-money straddle with strike close to the relevant forward price,
use the shorthand that any required discounting, dividend, and spot--forward
conversions have already been made, and define the observed premium on the
event-window return scale as

\[
\boxed{
m_{IV}\equiv\frac{C_{ATM}+P_{ATM}}{P_0}
\approx\E^{\mathbb Q}|\Rp|.
}
\]

The approximation suppresses the discounting, dividend, spot--forward, and
strike adjustments needed to express the straddle premium as a return moment.
Those adjustments must be made explicitly in an empirical implementation. It
also does not identify the physical distribution. Under a centered short-horizon
normal reference model,

\[
m_{IV}\approx\sqrt{\frac{2}{\pi}}\,\sigma\sqrt{T}
\]

\newpage

Here $\sigma$ is the reference volatility scale per unit time and $T$ is the
event-horizon length measured in the same time units. The normal expression is a
derived identity conditional on that reference distribution; it is not a general
property of the straddle. The risk-neutral
absolute move is not automatically a physical expected move. Let $h>0$ denote
an explicit physical-to-risk-neutral scale adjustment. The value and even the
direction of that adjustment are not identified by the theory and no numerical
value is assigned to it here.

\subsection{Tail magnitude in a symmetric three-state model}

If $\Rp\in\{-m,0,+m\}$, then

\[
\E|\Rp|=m(p_B+p_U).
\]

If the neutral state were a point mass at zero and $p_B+p_U>0$, consistency
with the adjusted absolute-move budget $h m_{IV}$ would require

\[
m=\frac{h m_{IV}}{p_B+p_U}.
\]

This is already different from setting $m=m_{IV}$: that shortcut ignores both
tail concentration and the unidentified physical-to-risk-neutral scale
adjustment, so it is not generally justified.

\subsection{Neutral-band leakage under a normal reference}
\label{sec:leak}

The neutral state is an interval in the underlying return distribution, even
though the three-state payoff approximation represents it by one value. To
derive a reference correction, suppose a centered normal return
$X\sim\mathcal N(0,s_{\mathcal N}^2)$ is used and the neutral band is
$|X|\le z s_{\mathcal N}$, so the threshold in Section~\ref{sec:setup} is
$a=z s_{\mathcal N}$. Its reference probability is
$p_N^{\mathrm{ref}}=2\Phi(z)-1$, and

\[
\E|X|=2s_{\mathcal N}\varphi(0),
\]

while the absolute move inside the neutral band is

\[
\E\left[|X|\mathbf 1\{|X|\le z s_{\mathcal N}\}\right]
=2s_{\mathcal N}\left(\varphi(0)-\varphi(z)\right).
\]

The fraction of the absolute-move budget absorbed by the neutral state is

\[
\boxed{
\ell_{\mathcal N}(p_N^{\mathrm{ref}})=1-\frac{\varphi(z)}{\varphi(0)}
=1-e^{-z^2/2},
\qquad
z=\Phi^{-1}\left(\frac{1+p_N^{\mathrm{ref}}}{2}\right).
}
\]

If this normal reference is also used to construct the physical tail budget,
define

\[
\boxed{
M=h m_{IV}\bigl(1-\ell_{\mathcal N}(p_N^{\mathrm{ref}})\bigr).
}
\]

If a conditional model supplies a physical $p_N$, that probability must remain
separate from $p_N^{\mathrm{ref}}$; replacing the normal-reference probability
with a conditional estimate would be an additional assumption. In the general
model, $M$ is therefore a separately specified or estimated physical tail
absolute-move budget. Given that budget, the symmetric three-state magnitude is

\[
\boxed{
m=\frac{M}{p_B+p_U}.
}
\]

\subsection{A reference distributional mapping, not an options-implied distribution}

The straddle identifies an absolute-move moment under the risk-neutral measure;
it does not identify physical probabilities for the three states. To obtain
state probabilities, one must specify a physical reference distribution and a
threshold $a$:

\[
\boxed{
\begin{aligned}
p_B^{\mathrm{ref}}&=\Prob_{\mathbb P}(\Rp<-a),\\
p_N^{\mathrm{ref}}&=\Prob_{\mathbb P}(|\Rp|\le a),\\
p_U^{\mathrm{ref}}&=\Prob_{\mathbb P}(\Rp>a).
\end{aligned}}
\]

For example, if a centered normal physical reference with scale
$s_{\mathbb P}>0$ is imposed, then

\[
p_N^{\mathrm{ref}}=2\Phi\left(\frac{a}{s_{\mathbb P}}\right)-1,
\qquad
p_B^{\mathrm{ref}}=p_U^{\mathrm{ref}}
=\frac{1-p_N^{\mathrm{ref}}}{2}.
\]

Equating $s_{\mathbb P}$ with a risk-neutral scale inferred from $m_{IV}$ would
be an additional physical/risk-neutral equivalence assumption, not a consequence
of the straddle. The paper therefore does not assign numerical values to a
supposed ``options-implied'' probability vector. It uses only the explicitly
labelled symmetry benchmark $p_B=p_U$ with $p_N$ left symbolic. With symmetric
tail magnitudes, that benchmark implies

\[
\boxed{
EV_0=m(p_U-p_B)=0
}
\]

by identity. A zero produced by this benchmark is not evidence that a predictive
model reconstructed market efficiency; no predictive model has been evaluated.

%======================================================================
\section{Scenario Probabilities}
\label{sec:prob}
%======================================================================

\subsection{State labels and probability sources}

The thresholds defining $B$, $N$, and $U$ must be fixed before outcomes are
observed. Once they are fixed, a probability vector can come from one of three
distinct sources:

\begin{enumerate}
\item the market-implied absolute-move scale together with a separately stated
      physical probability benchmark;
\item an unconditional historical frequency;
\item a fitted conditional model $\hat p_i(x)=\Prob(Y=i\mid X=x)$.
\end{enumerate}

These sources must not be reported as though they were interchangeable. In
particular, the absolute-move moment from options does not supply the physical
state probabilities, and any benchmark vector is an assumption rather than a
result of the conditional model.

\subsection{Ordered probability model}

Because $B\prec N\prec U$, one possible parametric specification is an
ordered-logit model. This is a modeling choice, not a consequence of the state
ordering. Let

\[
y^*=x^{\!\top}\beta+\varepsilon,
\qquad
\varepsilon\sim\operatorname{Logistic}(0,1),
\qquad c_1<c_2.
\]

The unit scale in the logistic error is a normalization for identification, not
a calibrated dispersion value. With $\Lambda(z)=1/(1+e^{-z})$,

\[
\boxed{
\begin{aligned}
\Prob(Y=B\mid x)&=\Lambda(c_1-x^{\!\top}\beta),\\
\Prob(Y=N\mid x)&=\Lambda(c_2-x^{\!\top}\beta)-\Lambda(c_1-x^{\!\top}\beta),\\
\Prob(Y=U\mid x)&=1-\Lambda(c_2-x^{\!\top}\beta).
\end{aligned}}
\]

The ordering reduces the parameter count and imposes the economically coherent
restriction that a feature shifting probability toward $U$ shifts it away from
$B$. An unrestricted multinomial model can be used, but it requires an explicit
reference-class constraint; otherwise its latent scores are unidentified under
common shifts.

\subsection{Separate surprise and reaction channels}

Fundamental variables may predict the earnings surprise, while positioning
variables may predict the price reaction conditional on that surprise. A
conceptual factorization is

\[
\boxed{
\Prob(\Rp)
=\sum_s \Prob(\Rp\mid s,\text{positioning})
\Prob(s\mid\text{fundamentals}).
}
\]

This distinction is structural. A strong earnings result need not produce a
positive event return if optimistic expectations were already priced. Combining
the channels without a stated factorization can turn a reaction effect into a
spurious earnings-quality effect.

\subsection{Regularization and temporal ordering}

If the ordered-logit model is estimated, one possible regularized objective is

\[
\hat\beta
=\arg\min_\beta\left[
-\sum_{t=1}^{T}\ln\Prob(Y_t\mid X_t;\beta)
 +\eta\|\beta\|_2^2
\right],
\]

where $\eta\ge 0$ is selected using time-ordered validation. Standardization
constants must also be computed on the training window alone. These are
requirements for a future empirical study, not evidence that the model has
already been fitted.

%======================================================================
\section{From Fair Value to Event Price}
\label{sec:kappa}
%======================================================================

\subsection{The convergence coefficient}

Fair value and market price are different variables. A minimal bridge is

\[
\boxed{
P_1=P_0+\kappa(V_1-P_0)+\epsilon,
\qquad
\E[\epsilon\mid V_1]=0.
}
\]

Dividing by $P_0$ gives

\[
\boxed{
\Rp=\kappa\Rv+\frac{\epsilon}{P_0},
\qquad
\E[\Rp]=\kappa\E[\Rv].
}
\]

The coefficient $\kappa$ is not a convenience parameter. It is the statement
about how much of a modeled value gap becomes price movement over the event
horizon. The core theory leaves its sign and magnitude unrestricted. An
application may impose a partial-convergence or no-overshoot restriction, but
that would be an additional modeling assumption. Treating the entire modeled
gap as an immediate event return is the special full-convergence case and is not
implied by the valuation equations.

\subsection{Why an intercept is required}

An empirical version of the bridge should be written across events as

\[
\Rp_{j,t}=\alpha+\kappa\Rv_{j,t}+\epsilon_{j,t}.
\]

The intercept is not optional. A common event premium, omitted variable, or
other non-valuation component belongs in $\alpha$ rather than being forced into
$\kappa$. A regression through the origin can therefore attribute a mean event
return to the valuation gap and inflate the apparent convergence coefficient.

\subsection{Measurement error}

The valuation gap is itself estimated. If the measured regressor equals
$\Rv+u$, where $u$ is mean-zero classical measurement error independent of the
true regressor and structural error with variance $\varsigma_V^2/P_0^2$, then
the usual attenuation expression is

\[
\boxed{
\operatorname{plim}\hat\kappa
=\kappa\frac{\Var(\Rv)}
{\Var(\Rv)+\varsigma_V^2/P_0^2}.
}
\]

The direction is toward zero. A future empirical study must therefore report
the valuation uncertainty and must not interpret an uncorrected insignificant
slope as proof that valuation has no relation to event returns.

\subsection{No double counting}

There are two related but distinct ways to describe an event return:

\begin{enumerate}
\item a scenario-payoff model assigns $\Rp_i$ directly to reaction states and
      computes $\sum_i p_i\Rp_i$;
\item a valuation model assigns $\Rv_i$ and maps it into price through
      $\kappa\sum_i p_i\Rv_i$.
\end{enumerate}

If the valuation gap is already used to define the state payoffs, the second
quantity is not an additional return to be added on top. A complete empirical
model must specify whether the valuation channel changes the state magnitudes,
the state probabilities, or a separate component of the price innovation.

%======================================================================
\section{Expected Value and Closed Forms}
\label{sec:ev}
%======================================================================

\subsection{General expected event return}

For state-dependent price returns,

\[
\boxed{
EV=\E[\Rp]=\sum_{i\in\{B,N,U\}}p_i\Rp_i.
}
\]

The notation $EV$ means expected value in the statistical sense only. Enterprise
value is $\Theta$ throughout.

For the symmetric event-payoff model $\Rp\in\{-m,0,+m\}$,

\[
\boxed{
EV=m(p_U-p_B).
}
\]

The directional edge is therefore the probability tilt $p_U-p_B$ multiplied by
the event magnitude. A separately imposed symmetric benchmark has zero tilt by
construction; the straddle alone does not impose that symmetry.

\subsection{Asymmetric event payoffs}
\label{sec:asym}

Let

\[
\Rp\in\{-m_B,0,+m_U\},
\qquad
\rho=\frac{m_B}{m_U}>0.
\]

The adjusted tail absolute-move budget $M$ supplies the moment equation

\[
p_Bm_B+p_Um_U=M.
\]

Solving with $m_B=\rho m_U$ gives

\[
\boxed{
m_U=\frac{M}{\rho p_B+p_U},
\qquad
m_B=\frac{\rho M}{\rho p_B+p_U}.
}
\]

Expected value is then

\[
\boxed{
EV=M\frac{p_U-\rho p_B}{\rho p_B+p_U}.
}
\]

The zero-edge condition is

\[
\boxed{
EV=0\quad\Longleftrightarrow\quad p_U=\rho p_B.
}
\]

The symmetric condition $p_U=p_B$ is a special case, not a general result.

\subsection{Symmetric closed forms}

For the symmetric model, write, for a nondegenerate tail and return variance,

\[
t=p_U-p_B,
\qquad
q=p_U+p_B,
\qquad
m=\frac{M}{q}.
\]

The mean, variance, risk-adjusted ratio, and exact discrete Kelly fraction are

\[
\boxed{
\begin{aligned}
EV&=\frac{Mt}{q},\\[3pt]
\Var(\Rp)&=\left(\frac{M}{q}\right)^2(q-t^2),\\[3pt]
\frac{EV}{\sigma_{\Rp}}&=\frac{t}{\sqrt{q-t^2}},\\[3pt]
f^*&=\frac{t}{M}.
\end{aligned}}
\]

The risk-adjusted ratio is independent of the absolute move scale. The exact
Kelly expression follows by maximizing
$\sum_i p_i\ln(1+f\Rp_i)$ over the three outcomes; it is not a recommendation
to use Kelly sizing when probabilities are uncertain.

\subsection{Neutral probability and structural invariance}

Holding the tail ratio $p_U:p_B$ fixed, the neutral probability affects the
event mean only through the leakage term in $M$ when $M$ is constructed from the
normal reference in Section~\ref{sec:leak}. If $M$ is separately estimated, the
neutral probability does not enter the expected-value formula in this
three-state payoff specification. This is a useful diagnostic: a complicated
three-state model does not automatically gain decision information from
estimating $p_N$.

%======================================================================
\section{Risk and Estimation Error}
\label{sec:error}
%======================================================================

\subsection{Downside risk}

Variance penalizes positive and negative movement symmetrically. For a threshold
$\tau$ and power $n\ge 1$, define the lower partial moment

\[
\boxed{
\operatorname{LPM}_n(\tau)
=\sum_i p_i\left[\max(\tau-\Rp_i,0)\right]^n.
}
\]

For the conventional zero-return benchmark $\tau=0$ (a choice of benchmark,
not a theoretical implication), downside deviation is
$\sqrt{\operatorname{LPM}_2(0)}$. In the symmetric three-state model,

\[
\operatorname{LPM}_1(0)=p_Bm,
\qquad
\sqrt{\operatorname{LPM}_2(0)}=m\sqrt{p_B}.
\]

The resulting Sortino-like ratio is $t/\sqrt{p_B}$, so it is also degenerate
under symmetry. A downside-sensitive statistic becomes informative only after
asymmetry is admitted.

\subsection{Sampling-error floor}

For unconditional multinomial proportions based on independent event
observations with a fixed probability vector,

\[
\boxed{
\Var(\hat p_U-\hat p_B)
=\frac{p_U(1-p_U)+p_B(1-p_B)+2p_Up_B}{n}.
}
\]

Therefore, in the symmetric model, the corresponding multinomial sampling
floor for the expected return is

\[
\boxed{
\se_{\mathrm{floor}}(\widehat{EV})
=m\sqrt{\frac{p_U(1-p_U)+p_B(1-p_B)+2p_Up_B}{n}}.
}
\]

This is a lower bound on the uncertainty of a fitted conditional model, not its
standard error. A conditional probability estimate also carries coefficient
uncertainty, covariance, model-selection uncertainty, uncertainty in $m$, and
possible dependence across event observations. A model that reports a point
estimate without an interval is therefore incomplete.

\subsection{Scale-free significance}

Holding $m$ fixed in the multinomial sampling calculation, both $EV$ and its
sampling floor scale linearly in $m$, so their ratio does not depend on the
absolute event magnitude. Changing the straddle-to-tail
mapping can change the dollar or percentage estimate without changing the basic
significance problem. Directional information is the bottleneck.

\subsection{Illustrative sensitivity}

Let $(\tilde p_B,\tilde p_N,\tilde p_U)$ denote any hypothetical conditional
probability vector with $\tilde p_U>\tilde p_B$. It is distinct from the
symmetric benchmark and is used only to show how a directional tilt enters the
formulas. Under symmetry,

\[
EV_{\mathrm{sym}}
=M\frac{\tilde p_U-\tilde p_B}{\tilde p_B+\tilde p_U}.
\]

If the downside magnitude is asymmetric, then

\[
EV_{\mathrm{asym}}
=M\frac{\tilde p_U-\rho\tilde p_B}
{\rho\tilde p_B+\tilde p_U}.
\]

The edge is exactly zero at $\rho=\tilde p_U/\tilde p_B$ and changes sign
above that value. No numerical probability vector or physical-move haircut is
needed to establish this result.

%======================================================================
\section{Theoretical Implications and Empirical Program}
\label{sec:implications}
%======================================================================

\subsection{What the theory predicts}

The framework produces testable implications without claiming that they have
already been tested:

\begin{enumerate}
\item A fitted probability model must beat both an unconditional historical
      base rate and a pre-specified physical benchmark on proper out-of-sample
      scoring rules if it is to claim incremental probability information.
\item The valuation layer should be evaluated through an intercept-inclusive
      estimate of $\kappa$, with measurement error in the value gap reported.
\item The downside-to-upside ratio $\rho$ should be estimated before a
      directional expected value is interpreted, because it can reverse the
      sign of the result.
\item Any economic evaluation should use event returns and appropriate
      transaction-cost, liquidity, and alternative-investment benchmarks rather
      than treating zero return as the only alternative.
\end{enumerate}

\subsection{A falsification protocol}

A future empirical study should use time-ordered training and test periods,
walk-forward refitting, and probability calibration metrics such as Brier score,
log loss, and expected calibration error. For probabilities $\hat p_{i,t}$ and
realized state $Y_t$,

\[
\operatorname{Brier}
=\frac1T\sum_{t=1}^{T}\sum_i
\left(\hat p_{i,t}-\mathbf 1\{Y_t=i\}\right)^2,
\qquad
\operatorname{LogLoss}
=-\frac1T\sum_{t=1}^{T}\ln\hat p_{Y_t,t}.
\]

The economic test is separate from the probability test. A well-calibrated
probability model need not generate a profitable strategy after costs, and a
profitable strategy need not identify which part of the return came from
valuation, positioning, or a broad announcement premium.

\subsection{Position in the literature}

The framework is consistent with the literature's distinction between event
information, announcement returns, and post-earnings drift. Earnings
announcements have long been treated as information events
\cite{beaver1968,ball1968,chari1988}. Cross-sectional announcement premia and
their limits to arbitrage have been documented in diversified portfolios
\cite{cohen2007,frazzini2007,savor2016,barber2013}. Post-earnings drift is a
separate phenomenon because the surprise is known after the event
\cite{bernard1989,linnainmaa2019}.

Those findings motivate the model's caution; they do not validate a particular
single-name forecast or any parameter used in the illustrations.

%======================================================================
\section{Limitations and Boundary Conditions}
\label{sec:limitations}
%======================================================================

\begin{enumerate}
\item The three-state approximation discards fat tails and path dependence.
      This is a deliberate low-dimensional representation, not a claim that
      the true event distribution has three points.
\item $h$, the physical adjustment from the risk-neutral straddle moment, is
      not identified here.
\item $\rho$ is not estimated. Since it determines the zero-edge condition, no
      asymmetric expected-value conclusion is empirical in this paper.
\item $\kappa$ is not estimated. A fair-value gap does not become an event
      return without an empirically supported price-convergence relation.
\item The surprise variables require forecasts of actual earnings, revenue,
      guidance, and persistence. If those forecasts add no information beyond
      consensus, the fundamental branch contributes no incremental edge.
\item Implied-volatility crush, liquidity, trading halts, guidance withdrawal,
      taxes, and transaction costs require separate models when an economic
      application is attempted.
\item The straddle identifies an absolute-move moment, not physical state
      probabilities. Any numerical probability benchmark is an explicit
      distributional assumption and should not be described as a forecast.
\end{enumerate}

%======================================================================
\section{Conclusion}
%======================================================================

The main theoretical result is a negative one in the useful scientific sense:
a pre-earnings expected-value claim is underdetermined when it reports only a
fair-value gap and a symmetric implied move. The straddle supplies an absolute
move moment, not a tail magnitude. A neutral band absorbs part of that moment.
Directional value depends on a probability tilt, while the sign under realistic
asymmetry depends on the ratio of downside to upside magnitudes. A valuation gap
reaches price only through a convergence coefficient, and the uncertainty of the
probability tilt must be quantified relative to the event counts available for a
single name.

The framework therefore defines a disciplined research object rather than a
validated trading system. Its null is explicit, its parameters are separated,
its assumptions are visible, and its empirical requirements are stated. A future
study may estimate the parameters and test the implications. Until then, the
theory's appropriate conclusion is not that an edge exists, but that any claimed
edge must survive the magnitude, asymmetry, convergence, and uncertainty tests
derived here.

\appendix

%======================================================================
\section{Secondary Decision Quantities}
\label{app:decision}
%======================================================================

This appendix records quantities that are useful for an eventual application but
are not needed for the theoretical contribution.

\subsection{Mean-variance utility}

For an action $a$ with event return $\Rp_a$, a local mean-variance approximation
to a CRRA objective, after normalizing by a reference wealth level, is

\[
U(a)=\E[\Rp_a]-\frac{\gamma}{2}\Var(\Rp_a)-C(a),
\]

where $\gamma>0$ is an unspecified relative-risk-aversion parameter and $C(a)$
collects trading, tax, liquidity, and other action-specific costs. The
comparison is against the best alternative, not automatically against zero. For
the costless analytical benchmark with the alternative normalized to zero, the
break-even risk-aversion coefficient in the symmetric model is

\[
\boxed{
\gamma^*=\frac{2tq}{M(q-t^2)}.
}
\]

This expression is a sensitivity result. It is not a recommendation to infer an
investor's risk aversion from a single event.

\subsection{Kelly sizing}

Under known probabilities, no transaction costs, and no position or leverage
constraints, the exact growth-optimal fraction in the symmetric three-state
model, over the feasible domain $1+f\Rp_i>0$, is

\[
f^*=\frac{t}{M}.
\]

The exact expression does not account for parameter uncertainty, costs,
constraints, or model risk. Any practical use would therefore require an
externally specified decision rule for shrinking or otherwise restricting the
fraction.

%======================================================================
\section{Illustrative Calculation}
\label{app:illustration}
%======================================================================

The following calculation is included only to make the identities concrete. It
is not a stock forecast and uses no observed event data.

Let $a$ be a pre-registered neutral threshold and let a hypothetical centered
normal physical reference distribution be $\mathcal N(0,s_{\mathbb P}^2)$, with
$s_{\mathbb P}>0$. Then

\[
p_N^{\mathrm{ref}}=2\Phi\left(\frac{a}{s_{\mathbb P}}\right)-1,
\qquad
p_B^{\mathrm{ref}}=p_U^{\mathrm{ref}}
=\frac{1-p_N^{\mathrm{ref}}}{2}.
\]

For any hypothetical conditional vector
$(\tilde p_B,\tilde p_N,\tilde p_U)$ with nonnegative components summing to
one, define

\[
\tilde q=\tilde p_B+\tilde p_U,
\qquad
\tilde t=\tilde p_U-\tilde p_B,
\qquad
M>0.
\]

Here $M$ is a separately specified physical tail absolute-move budget; it is
not inferred from $\tilde p_N$ unless the additional normal-reference mapping
and physical-scale assumptions are imposed.

The symmetric and asymmetric expected values are respectively

\[
EV_{\mathrm{sym}}=M\frac{\tilde t}{\tilde q},
\qquad
EV_{\mathrm{asym}}=M\frac{\tilde p_U-\rho\tilde p_B}
{\rho\tilde p_B+\tilde p_U}.
\]

When $\tilde p_B>0$, the asymmetric edge changes sign at
$\rho=\tilde p_U/\tilde p_B$. This symbolic result is the point of the
illustration; no arbitrary numerical probability vector is needed.

%======================================================================
\addcontentsline{toc}{section}{References}
\begin{thebibliography}{99}

\bibitem{ball1968}
Ball, R. and Brown, P. (1968). An Empirical Evaluation of Accounting Income
Numbers. \emph{Journal of Accounting Research}, 6(2), 159--178.

\bibitem{barber2013}
Barber, B. M., De George, E. T., Lehavy, R. and Trueman, B. (2013). The
Earnings Announcement Premium Around the Globe. \emph{Journal of Financial
Economics}, 108(1), 118--138.

\bibitem{beaver1968}
Beaver, W. H. (1968). The Information Content of Annual Earnings Announcements.
\emph{Journal of Accounting Research}, 6, 67--92.

\bibitem{bernard1989}
Bernard, V. L. and Thomas, J. K. (1989). Post-Earnings-Announcement Drift:
Delayed Price Response or Risk Premium? \emph{Journal of Accounting Research},
27, 1--36.

\bibitem{chari1988}
Chari, V. V., Jagannathan, R. and Ofer, A. R. (1988). Seasonalities in Security
Returns: The Case of Earnings Announcements. \emph{Journal of Financial
Economics}, 21(1), 101--121.

\bibitem{cohen2007}
Cohen, D. A., Dey, A., Lys, T. Z. and Sunder, S. V. (2007). Earnings
Announcement Premia and the Limits to Arbitrage. \emph{Journal of Accounting
and Economics}, 43(2--3), 153--180.

\bibitem{frazzini2007}
Frazzini, A. and Lamont, O. (2007). The Earnings Announcement Premium and
Trading Volume. NBER Working Paper No. 13090.

\bibitem{linnainmaa2019}
Linnainmaa, J. T. and Zhang, C. Y. (2019). The Earnings Announcement Return
Cycle. Working paper, SSRN 3183318.

\bibitem{savor2016}
Savor, P. and Wilson, M. (2016). Earnings Announcements and Systematic Risk.
\emph{Journal of Finance}, 71(1), 83--138.

\end{thebibliography}

\end{document}
